%% OmniLaTeX Color Theme Preview
%%
%% Demonstrates all 6 built-in themes with dark/light mode
%% and institution color integration.
%%
%% Build:  lualatex main.tex

\documentclass[
    language=english,
    doctype=article,
    institution=none,
    titlestyle=simple,
]{../../omnilatex}

\RequirePackage{lib/utils/omnilatex-themes}
\RequirePackage{config/document-settings}

\title{OmniLaTeX Color Themes}
\subtitle{A visual comparison of all built-in color themes}
\author{OmniLaTeX}
\date{\today}

\begin{document}

\maketitle

\tableofcontents
\newpage

%% ── Theme Comparison Table ────────────────────────

\section{Theme Overview}

\begin{table}[htbp]
    \centering
    \begin{tabular}{lll}
        \textbf{Theme} & \textbf{Character} & \textbf{Dark Variant} \\
        \hline
        default    & White bg, blue accents   & monochrome-dark (fallback) \\
        midnight   & Navy bg, cyan accents    & built-in (always dark) \\
        forest     & Dark green, green accents & built-in (always dark) \\
        rose       & Dark pink, magenta accents & built-in (always dark) \\
        monochrome & Black/white only         & separate dark palette \\
        sepia      & Warm paper tones         & monochrome-dark (fallback) \\
    \end{tabular}
    \caption{Available color themes and their dark mode behavior.}
\end{table}

\newpage

%% ── Default Theme ─────────────────────────────────

\section{Default Theme}

\usetheme{default}

This is the \textbf{default} OmniLaTeX theme: a clean white background with
dark text and blue accents. It is designed for maximum readability in print
and on screen.

\subsection{Elements}

\begin{itemize}
    \item Body text in standard weight
    \item \textbf{Bold text} for emphasis
    \item \textit{Italic text} for secondary emphasis
    \item A \href{https://github.com/WyattAu/OmniLaTeX-template}{hyperlink} in accent color
\end{itemize}

\begin{block}{Block Environment}
    Blocks use the accent color for their title bar and a tinted background.
\end{block}

\begin{addmargin}{1em}
\colorbox{red!10}{\parbox{0.95\linewidth}{\textbf{Alert Block}\\Alert blocks retain their own red color regardless of theme.}}
\end{addmargin}

\begin{addmargin}{1em}
\colorbox{green!10}{\parbox{0.95\linewidth}{\textbf{Example Block}\\Example blocks retain their green color regardless of theme.}}
\end{addmargin}

\newpage

%% ── Midnight Theme ────────────────────────────────

\section{Midnight Theme}

\usetheme{midnight}

This is the \textbf{midnight} theme: a dark navy background with cyan accents.
It is inherently dark — there is no light variant.

\subsection{Elements}

\begin{itemize}
    \item Body text in light weight on dark background
    \item \textbf{Bold text} for emphasis
    \item \textit{Italic text} for secondary emphasis
    \item A \href{https://github.com/WyattAu/OmniLaTeX-template}{hyperlink} in cyan
\end{itemize}

\begin{block}{Block Environment}
    Blocks use cyan-tinted title bars against the dark background.
\end{block}

\begin{addmargin}{1em}
\colorbox{red!10}{\parbox{0.95\linewidth}{\textbf{Alert Block}\\Alert blocks retain their red color on the dark background.}}
\end{addmargin}

\begin{addmargin}{1em}
\colorbox{green!10}{\parbox{0.95\linewidth}{\textbf{Example Block}\\Example blocks retain their green color on the dark background.}}
\end{addmargin}

\newpage

%% ── Forest Theme ──────────────────────────────────

\section{Forest Theme}

\usetheme{forest}

This is the \textbf{forest} theme: a dark green background with green accents.
Like midnight, it is inherently dark.

\subsection{Elements}

\begin{itemize}
    \item Body text in light weight on dark green background
    \item \textbf{Bold text} for emphasis
    \item \textit{Italic text} for secondary emphasis
    \item A \href{https://github.com/WyattAu/OmniLaTeX-template}{hyperlink} in green
\end{itemize}

\begin{block}{Block Environment}
    Blocks use green-tinted title bars against the dark background.
\end{block}

\newpage

%% ── Rose Theme ────────────────────────────────────

\section{Rose Theme}

\usetheme{rose}

This is the \textbf{rose} theme: a dark pink background with magenta accents.
It is inherently dark.

\subsection{Elements}

\begin{itemize}
    \item Body text in light weight on dark pink background
    \item \textbf{Bold text} for emphasis
    \item \textit{Italic text} for secondary emphasis
    \item A \href{https://github.com/WyattAu/OmniLaTeX-template}{hyperlink} in magenta
\end{itemize}

\begin{block}{Block Environment}
    Blocks use magenta-tinted title bars against the dark background.
\end{block}

\newpage

%% ── Monochrome Theme ──────────────────────────────

\section{Monochrome Theme (Light)}

\usetheme{monochrome}

This is the \textbf{monochrome} theme in light mode: pure black and white
with no color whatsoever.

\subsection{Elements}

\begin{itemize}
    \item Body text in standard black
    \item \textbf{Bold text} for emphasis
    \item \textit{Italic text} for secondary emphasis
    \item A \href{https://github.com/WyattAu/OmniLaTeX-template}{hyperlink} in dark gray
\end{itemize}

\begin{block}{Block Environment}
    Blocks use grayscale title bars.
\end{block}

\newpage

%% ── Monochrome Theme (Dark) ───────────────────────

\section{Monochrome Theme (Dark)}

\usetheme{monochrome-dark}

This is the \textbf{monochrome} theme in dark mode: white text on a near-black
background with no color.

\subsection{Elements}

\begin{itemize}
    \item Body text in off-white
    \item \textbf{Bold text} for emphasis
    \item \textit{Italic text} for secondary emphasis
    \item A \href{https://github.com/WyattAu/OmniLaTeX-template}{hyperlink} in light gray
\end{itemize}

\begin{block}{Block Environment}
    Blocks use grayscale title bars on dark background.
\end{block}

\newpage

%% ── Sepia Theme ───────────────────────────────────

\section{Sepia Theme}

\usetheme{sepia}

This is the \textbf{sepia} theme: warm, paper-like tones designed for
comfortable reading. It evokes the look of aged paper.

\subsection{Elements}

\begin{itemize}
    \item Body text in warm dark brown
    \item \textbf{Bold text} for emphasis
    \item \textit{Italic text} for secondary emphasis
    \item A \href{https://github.com/WyattAu/OmniLaTeX-template}{hyperlink} in warm orange-brown
\end{itemize}

\begin{block}{Block Environment}
    Blocks use warm-toned title bars on the sepia background.
\end{block}

\newpage

%% ── Dark/Light Toggle ─────────────────────────────

\section{Dark/Light Mode Toggle}

The theme system provides \texttt{\textbackslash darkmode} and
\texttt{\textbackslash lightmode} commands to toggle between modes
mid-document.

\subsection{Starting in Light Mode}

\usetheme{default}

This paragraph is rendered in the \textbf{default} (light) theme.
Links are \href{https://github.com/WyattAu/OmniLaTeX-template}{blue}.

\subsection{Switching to Dark Mode}

\darkmode

After calling \texttt{\textbackslash darkmode}, the page switches to the
monochrome-dark palette (the default theme has no explicit dark variant,
so it falls back to monochrome-dark).

Links are now \href{https://github.com/WyattAu/OmniLaTeX-template}{light gray}.

\subsection{Switching Back to Light Mode}

\lightmode

Calling \texttt{\textbackslash lightmode} restores the original theme.
Links are \href{https://github.com/WyattAu/OmniLaTeX-template}{blue} again.

\newpage

%% ── Institution Color Integration ─────────────────

\section{Institution Color Integration}

\usetheme{midnight}

When an institution is configured (e.g., \texttt{institution=tuhh}), the
institution's \texttt{universityprimary} color overrides the theme's accent,
link, and block frame colors. This ensures institutional branding is preserved
across all themes.

\subsection{How It Works}

\begin{enumerate}
    \item The institution config defines \texttt{universityprimary} and
          \texttt{universitysecondary} colors.
    \item When \texttt{\textbackslash usetheme} is called, it checks whether
          \texttt{universityprimary} is defined.
    \item If present, the institution color replaces the theme's accent,
          link, and block frame colors.
    \item The background, foreground, and other semantic colors remain
          governed by the theme.
\end{enumerate}

\begin{block}{Custom Override}
    You can also use \texttt{\textbackslash setthemecolor\{accent\}\{red\}}
    to override individual color slots after applying a theme.
\end{block}

\newpage

%% ── Custom Color Override ─────────────────────────

\section{Custom Color Overrides}

\usetheme{default}

The \texttt{\textbackslash setthemecolor} command allows fine-grained
per-slot overrides. For example:

\begin{center}
    \texttt{\textbackslash setthemecolor\{accent\}\{red!70!black\}}\\[6pt]
    \texttt{\textbackslash setthemecolor\{link\}\{teal\}}\\[6pt]
    \texttt{\textbackslash setthemecolor\{heading\}\{purple\}}
\end{center}

Available slots: \texttt{bg}, \texttt{fg}, \texttt{heading}, \texttt{body},
\texttt{accent}, \texttt{blockbg}, \texttt{blockframe}, \texttt{link},
\texttt{rule}, \texttt{codebg}, \texttt{footernote}.

\subsection{Example Override}

\setthemecolor{accent}{red!70!black}

After overriding the accent color, this \textbf{bold text} still uses the
heading color, but footnote marks and caption labels now use red.

\end{document}
